\section{Related Work}\label{sec:related-work}

The research problem addressed by \AFD{} lies at the intersection of malware
taxonomy, secure messaging, adversarial file parsing, and content
reconstruction. These literatures use different units of analysis. Malware
research commonly studies a program's propagation or effects after execution.
Secure-messaging research studies confidentiality, authentication, and state
evolution across communicating endpoints. Parser-security research studies
the language recognized by software that consumes attacker-controlled bytes.
Content reconstruction studies whether an input can be transformed into a
narrower representation. A defensible attachment claim must not substitute the
conclusion of one literature for evidence required by another.

\subsection{Malware behavior and attachment carriers}

Cohen's foundational definition makes infection of other programs the
distinguishing operation of a computer virus~\cite{cohen1987}. Spafford's
analysis of the Internet Worm in turn illustrates propagation through network
services and trust relationships without requiring infection of a conventional
host file~\cite{spafford1989}. The distinction remains important because a
media attachment is not a virus merely because it contains suspicious bytes,
and a worm is not identified by the extension of the file through which an
operator first receives it. NIST's operational taxonomy separates viruses,
worms, Trojan horses, spyware, and other malicious code by behavior and
deployment context~\cite{nist80083}. Those categories provide threat context,
but they do not directly define a media grammar that a file-defense component
can recognize.

Behavioral ransomware defenses make this difference especially clear. UNVEIL
detects ransomware by observing interactions with user data and desktop state
in an instrumented environment~\cite{kharraz2016}. CryptoLock likewise studies
runtime indicators associated with destructive transformation of user
files~\cite{scaife2016}. Such systems can support a ransomware conclusion
because they observe behavior that instantiates the payload objective. \AFD{}
does not execute an attachment to obtain comparable evidence. It may block an
executable or scripted carrier that is labelled as ransomware in a corpus, but
the supported result is that the carrier violated the media-only policy. It is
not that ransomware behavior was inferred from the bytes.

This separation also applies to Trojans, staged-delivery roles, and spyware. A
Trojan is characterized by a deceptive relationship between represented and
actual function. A downloader obtains a follow-on payload from another source.
A dropper carries and places an embedded payload. A stager prepares or transfers
control to a follow-on stage. Spyware is characterized by covert collection or
disclosure. A filename, media type, or container property may contribute to
delivery of any of these programs, yet none establishes the downstream behavior
or role. The carrier taxonomy in \cref{sec:malware-taxonomy} therefore treats
malware labels and operational roles as threat and corpus strata while making
enforcement depend on observable file properties.

\subsection{Secure messaging and endpoint placement}

Secure-messaging analyses establish goals such as confidentiality, integrity,
authentication, forward secrecy, and recovery after compromise. The survey by
Unger et al. systematizes these properties and the design choices that support
them~\cite{unger2015}. Formal analysis of the Signal protocol provides stronger
reasoning about a particular messaging construction and its cryptographic
state transitions~\cite{cohngordon2020}. Neither line of work purports to prove
that plaintext obtained after successful decryption is safe for every local
parser. This is not a deficiency in those models. It is a separate endpoint
question.

The placement of \AFD{} follows from that separation. Outgoing reconstruction
occurs before encryption and incoming reconstruction occurs after decryption.
The relay can remain blind to attachment plaintext. The endpoint still needs a
local admissibility decision before publishing bytes to a renderer or another
application. The contribution is therefore complementary to end-to-end
encryption. It does not modify the cryptographic protocol or strengthen its
authentication theorem. It specifies what the endpoint may release after the
cryptographic layer has produced bytes.

Endpoint execution alone does not make two security designs equivalent.
Abelson et al. use \emph{client-side scanning} for systems that search endpoint
cleartext for centrally designated targets and may report a match or its source
to an outside party~\cite{abelson2024}. Their security and surveillance
analysis applies to that targeting and reporting architecture. \AFD{} is a
user-benefiting client-side defense. Its content-derived decision controls only
which reconstructed bytes the local application may receive, and it creates no
match report for the relay or a central authority. This distinction limits the
comparison to the shared fact that plaintext is processed on an endpoint.

\subsection{Parser differentials and language-oriented security}

Jana and Shmatikov demonstrated that discrepancies between the file-processing
logic of a detector and that of the protected host can enable evasion even when
the underlying executable behavior is unchanged~\cite{jana2012}. Their result
undermines any design that treats a leading signature, one MIME decision, or
one successful parser as an authoritative description of the whole object.
Language-theoretic security generalizes the problem by treating input handling
as language recognition and by emphasizing explicit, tractable input
languages~\cite{momot2016}.

Several systems improve the construction of trustworthy parsers. Nail derives
parsers and generators from a shared grammar and separates data description
from semantic processing~\cite{nail2014}. EverParse generates verified parsers
and serializers for binary message formats, with properties that include
memory safety, inverse parsing and serialization, and unique encodings in its
modeled domain~\cite{everparse2019}. Interval Parsing Grammars extend
declarative format descriptions to context-sensitive patterns such as
type--length--value fields and random-access references~\cite{ipg2023}. These
works motivate complete parsing, checked dependencies, and controlled
serialization. They do not by themselves choose an attachment policy, define
which media semantics may be retained, or establish safe publication in an
encrypted messenger.

The DARPA Safe Documents program is the closest program-level precedent for
combining format comprehension, safe subsetting, high-assurance translation,
and verified parser construction~\cite{darpaSafeDocs}. Its scope makes clear
that reducing ambiguous legacy languages and producing trustworthy parsers are
established research objectives. \AFD{} does not claim the verification
assurance pursued by SafeDocs. Its separate validator is an architectural
check over a policy-defined media profile, and it can still share a
specification or implementation error with the reconstruction handler.

\AFD{} adopts the language-oriented premise but makes a narrower engineering
composition. Name, supplied media type, and bounded byte evidence must converge
on a concrete format profile before dispatch under the evaluated policy. A
handler then performs a policy-required reconstruction level. A distinct
validator reparses the rebuilt output from offset zero to its exact terminal
boundary. Canonical component rules, resource bounds, output type agreement,
and post-reconstruction signatures are additional conjuncts. This composition
is an argument about one implementation and policy profile. It is not a proof
that all media decoders are correct.

\subsection{Content disarm and reconstruction}

Content disarm and reconstruction is direct prior art for transforming an
untrusted file into a narrower representation. Open Image CDR decodes JPEG
content to pixel state, rebuilds the image, and applies pixel manipulation to
address both non-image material and selected steganographic constructions while
measuring image usability~\cite{belkind2023}. ImSan, introduced by Koch et al.,
loads image content through Pillow and writes a new image with metadata and
extraneous file content removed to address image-based polyglots~\cite{koch2025}.
These systems establish semantic image reconstruction and image-polyglot
sanitization as precedents. \AFD{} therefore makes no priority claim for image
decode-and-reencode, metadata stripping, or image CDR.

PdfCDR provides the closest format-specific comparison outside raster media. It
defines and evaluates a PDF reconstruction process, including validation,
visual similarity, and conversion of detection rules into disarm
rules~\cite{dubin2023}. PDF remains an unsupported active class in the evaluated
\AFD{} media profile. The comparison is methodological rather than a claim that
\AFD{} supersedes PdfCDR for PDF. The residual question is how reconstruction
is selected, strengthened, independently validated, and transactionally
published at an encrypted-messaging endpoint under a pre-specified evidence
contract.

\subsection{Polyglots and semantic disagreement}

Polyglot files expose the limits of single-parser recognition. This manuscript
reserves \emph{true polyglot} for a whole octet string that is completely
accepted, from offset zero through the same terminal boundary, by at least two
incompatible parsers. A valid object followed by bytes that form or resemble
another object is instead an appended or trailing-object carrier unless both
parsers completely accept the whole string. Prefix camouflage and a nested
active object are likewise separate carrier constructions unless they meet the
complete dual-acceptance condition.

Koch et al. systematize the abuse and detection of polyglot files and evaluate
identification and sanitization techniques in that setting~\cite{koch2025}.
You et al. provide a focused demonstration of semantic gaps across ZIP parsers,
showing that nominal agreement on a format name does not guarantee agreement on
represented content~\cite{you2025}. The practical lesson is stronger than a
requirement to improve magic-byte detection. The accepted language needs exact
boundaries, parser-specific structure, and a policy for ambiguity.

The \AFD{} policy distinguishes true polyglots from prefix camouflage,
appended or trailing objects, and nested active objects. Parsing from offset
zero and exact terminal consumption address prefix and suffix constructions.
Independent incompatible complete parses establish a true-polyglot case.
Location-aware rules address nested active objects. A structural reconstruction
that copies encoded media cannot claim to eliminate a secondary interpretation
inside the copied region. A semantic reconstruction can support a stronger
neutralization statement only when it decodes to bounded semantic state,
copies no source container or compressed payload bytes, emits a new object,
and yields exactly one accepted output grammar. Otherwise the input is blocked.
Even then, information retained in pixels or samples remains outside the
harmlessness claim.

\begingroup
\footnotesize
\setlength{\tabcolsep}{3pt}
\begin{longtable}{L{0.18\textwidth}L{0.19\textwidth}L{0.26\textwidth}L{0.27\textwidth}}
\caption{Relationship between prior research objects and the claim made by
\AFD. A dash denotes a property outside the cited work's primary unit of
analysis, not a criticism of that work.}\label{tab:related-work-comparison}\\
\toprule
Research line & Primary observation & Result it can support & Boundary relative to \AFD{} \\
\midrule
\endfirsthead
\toprule
Research line & Primary observation & Result it can support & Boundary relative to \AFD{} \\
\midrule
\endhead
\midrule
\multicolumn{4}{r}{Continued on next page}\\
\endfoot
\bottomrule
\endlastfoot
Virus and worm theory~\cite{cohen1987,spafford1989} & Replication and propagation behavior & Behavioral distinction between infection and self-propagation after execution & Does not define canonical media acceptance \\
Ransomware behavior~\cite{kharraz2016,scaife2016} & Runtime modification of user files or desktop state & Detection of a payload objective in an observed execution & \AFD{} does not observe this behavior \\
Secure messaging~\cite{unger2015,cohngordon2020} & Protocol transcripts, keys, and compromise states & Communication security under stated assumptions & Decrypted attachment syntax is a separate endpoint problem \\
Reporting-oriented client-side scanning~\cite{abelson2024} & On-device search for designated targets with an external reporting path & Analysis of surveillance, security, and abuse risks & \AFD{} acts for the local user and sends no content-derived report \\
Detector differentials and LangSec~\cite{jana2012,momot2016} & Divergent interpretations and recognized languages & Evidence that parser agreement and explicit languages matter & Does not select a messaging publication policy \\
Generated or verified parsers~\cite{nail2014,everparse2019,ipg2023} & Grammar, parser, serializer, and proof obligations & Strong parser properties within a formalized domain & Does not establish semantic harmlessness of accepted media \\
DARPA SafeDocs~\cite{darpaSafeDocs} & Verified parsers, format comprehension, safe subsets, and high-assurance translators & Assurance for parsing and translation within defined languages & \AFD{} does not claim equivalent verification and adds an endpoint publication contract \\
Open Image CDR and ImSan~\cite{belkind2023,koch2025} & Semantic image rebuilding, metadata removal, and image-polyglot sanitization & Image-specific reconstruction and usability evidence & \AFD{} does not claim first image CDR and adds policy-indexed multi-format delivery conditions \\
PdfCDR~\cite{dubin2023} & PDF-specific validation, disarm rules, reconstruction, and visual comparison & Evidence for a PDF CDR process & PDF is blocked by the evaluated \AFD{} profile and no PDF superiority is claimed \\
Polyglot and parser-gap analysis~\cite{koch2025,you2025} & Complete whole-string acceptance by multiple formats or divergent meanings within a nominal format & Ambiguity classes and differential evidence & \AFD{} separates true polyglots from trailing, prefix, and nested carriers \\
\end{longtable}
\endgroup

The resulting novelty claim is deliberately compositional. \AFD{} is not a new
codec, general antivirus engine, behavioral sandbox, CDR primitive, or
proof-producing parser generator. Its residual novelty is endpoint-local,
user-benefiting reconstruction in an end-to-end encrypted workflow whose relay
remains blind. Exactly three routing sources select a profile, namely the
normalized portable name, caller-supplied media type, and bounded byte evidence.
An independent complete output parser then participates in delivery
authorization rather than routing. Policy-indexed reconstruction strength
prevents a weaker path from satisfying a stronger requirement. Independently
validated clean-only publication controls which bytes become deliverable. A
pre-specified evidence contract fixes the empirical obligations before results
are admitted. Each element has precedents, while their joint use defines the
claimed acceptance boundary.
